\chapter{Trabalhos Relacionados}
\label{chp:capitulo3}

Este capítulo revisa a literatura diretamente relacionada à predição de endpoint em processos siderúrgicos via aprendizado de máquina informado por física. A organização segue uma progressão do geral ao específico: parte dos fundamentos de PIML, avança para aplicações em fornos BOF e EAF, examina abordagens de geração de dados sintéticos industriais e, por fim, discute arquiteturas híbridas físico-informadas. A análise encerra com uma tabela comparativa e a identificação das lacunas que motivam a proposta desta dissertação.

%-------------------------------------------------------------------%
\section{Fundamentos de PIML Aplicados à Siderurgia}
\label{sec:piml-fundamentos}

A abordagem de Redes Neurais Informadas por Física (PINNs) foi formalizada por \citeauthor{raissi2019pinns} ao incorporar equações diferenciais parciais diretamente na função de perda do treinamento, permitindo que a rede neural resolva problemas diretos e inversos com supervisão parcial\,\cite{raissi2019pinns}. A formulação original demonstrou que, ao penalizar violações de leis físicas durante o gradiente descendente, é possível reduzir substancialmente a quantidade de dados rotulados necessária para treinar modelos precisos em domínios governados por equações de conservação. Esse paradigma abriu caminho para aplicações em transferência de calor, escoamento de fluidos e dinâmica estrutural.

Revisões abrangentes consolidaram o campo do PIML, estabelecendo vocabulário e taxonomia comuns\,\cite{cuomo2022physics,willard2020integrating,karniadakis2021physics}. A taxonomia adotada nesta dissertação, seguindo \citeauthor{willard2020integrating}, distingue três topologias de integração entre física e ML: \textit{serial} (um modelo físico alimenta, ou é realimentado por, um modelo de ML), \textit{paralela} (física e ML operam em ramos independentes cujas saídas são combinadas, tipicamente na função de perda, como nos PINNs clássicos\,\cite{raissi2019pinns}) e \textit{informada por atributos} (grandezas físicas pré-computadas entram como variáveis de entrada de um modelo de ML padrão, sem alterar seu objetivo de treinamento)\,\cite{willard2020integrating}. Esta última é a abordagem adotada nesta dissertação, na forma de engenharia de atributos pré-computados (Config~C, Capítulo~\ref{chp:capitulo4}). As três topologias são categorias idealizadas: sistemas reais frequentemente combinam elementos de mais de uma, e a fronteira entre elas é permeável dependendo do nível de abstração adotado na análise\,\cite{cuomo2022physics}. Um resultado transversal na literatura é que restrições físicas codificadas nos modelos aumentam a capacidade de extrapolação para regiões fora do domínio de treinamento --- propriedade especialmente relevante em processos industriais onde dados experimentais são escassos e caros\,\cite{willard2020integrating}.

%-------------------------------------------------------------------%
\section{Predição de Endpoint em Fornos Básicos a Oxigênio}
\label{sec:bof-endpoint}

A predição de endpoint em conversores BOF — temperatura final do banho, teor de carbono residual e, mais recentemente, fósforo — é um problema de regressão supervisionada com alta relevância industrial, pois erros de predição implicam reprocessamento dispendioso e desperdício energético\,\cite{chattopadhyay2022hybrid,zhou2024bof}. O uso de redes neurais para esse problema não é recente: \citeauthor{cox2002neural} já aplicavam computação neural à predição de endpoint em BOF há mais de duas décadas\,\cite{cox2002neural}, com desempenho e disponibilidade de dados de planta condicionados pelo estado da arte de ML da época. A literatura mais recente revisada nesta seção converge para arquiteturas que integram algum nível de conhecimento físico-químico, seja pela seleção cuidadosa de variáveis de entrada ou pela imposição de restrições no treinamento — o foco desta dissertação --- sem que isso implique que a modelagem estatística/ANN pré-PIML tenha sido superada de forma generalizada.

\citeauthor{xia2024dmpinn} propuseram o modelo dmPINNs (\textit{data-driven and mechanism-based Physics-Informed Neural Networks}) para predição do teor de carbono final em BOF\,\cite{xia2024dmpinn}. A arquitetura é \textit{paralela}: uma sub-rede extrai parâmetros latentes ($\alpha$, $\beta$) de uma equação cinética de descarburação apresentada pelos autores, que produz uma predição mecanicista $C_\text{cal}$; uma segunda sub-rede orientada a dados produz $C_\text{pre}$; o resultado final é a média ponderada $C_\text{dmPINN} = 0{,}5 \cdot C_\text{cal} + 0{,}5 \cdot C_\text{pre}$ (notação dos autores, distinta do $[\text{C}]_\text{final}$ desta dissertação), com função de perda que penaliza ambas as branches simultaneamente. Os resultados são reportados em termos de \textit{hit rates} para limiares de $\pm 0{,}009$, $\pm 0{,}012$ e $\pm 0{,}020$\,\%~em~massa, obtendo ganho de $+5{,}23$\,pp sobre o BPNN de linha de base no limiar mais restritivo. Esta arquitetura é fundamentalmente diferente da adotada nesta dissertação: o presente trabalho insere grandezas físicas como \textit{features} de entrada do modelo, não como restrição na função de perda.

\citeauthor{zhou2024bof} realizaram uma comparação sistemática entre algoritmos de ML e aprendizado profundo para predição de endpoint no BOF, com 13.528 corridas reais da Hesteel Group\,\cite{zhou2024bof}. Entre os modelos tabulares (45 variáveis brutas de composição, adições e O$_2$), CatBoost obteve o menor MAE para P ($2{,}28 \times 10^{-3}$\,\%~em~massa) e C ($6{,}92 \times 10^{-3}$\,\%~em~massa), enquanto LightGBM foi superior para temperatura (7{,}96\,°C). Os autores também incorporaram sete curvas de off-gas como séries temporais, combinadas ao modelo tabular via um \textit{backbone} TCN+ResCNN; esse modelo combinado reduz o MAE de teste dos três alvos simultaneamente --- não apenas de temperatura --- para $2{,}15 \times 10^{-3}$\,\% (P), $6{,}41 \times 10^{-3}$\,\% (C) e $5{,}78$\,°C (T), com MAE de $2{,}39 \times 10^{-3}$\,\%, $6{,}23 \times 10^{-3}$\,\% e $7{,}68$\,°C, respectivamente, em teste de campo real --- este último consistentemente pior que o de teste \textit{offline} para os três alvos, ilustrando o custo de generalização de laboratório para planta mesmo com dados reais. Nenhuma feature físico-derivada foi utilizada, o que confirma diretamente a lacuna endereçada por esta dissertação.

O trabalho de \citeauthor{chattopadhyay2022hybrid} apresentou um método híbrido para predição de endpoint em BOF com intercâmbio de informação serial em múltiplas etapas\,\cite{chattopadhyay2022hybrid}: modelos termodinâmicos de balanço de entalpia e oxigênio (MM\_T, MM\_C, MM\_P) fornecem estimativas iniciais de T, C e P, enquanto redes neurais independentes (ANN\_T, ANN\_C, ANN\_P) são treinadas nas variáveis brutas; na etapa híbrida, as predições das ANNs são realimentadas como entradas dos modelos teóricos — por exemplo, $\hat{C}_\text{ANN}$ substitui a assunção de carbono mediano no balanço entálpico de temperatura. Essa troca bidirecional reduz o NRMSE de predição de P em 9{,}77\% e de C em 3{,}7\% frente aos modelos teóricos isolados. A topologia é classificada como \textit{serial} na taxonomia de\,\cite{willard2020integrating} e constitui referência metodológica para a configuração C desta dissertação. \citeauthor{bae2020bof} demonstraram que grupos de atributos termodinâmicos elevam em 2–4 pontos percentuais as taxas de acerto para temperatura e fósforo, evidenciando que features físico-derivadas contribuem para estabilidade preditiva sob variações operacionais\,\cite{bae2020bof}.

A literatura em predição de endpoint BOF abrange desde regressão linear regularizada até redes neurais profundas\,\cite{zhou2024bof,chattopadhyay2022hybrid}. A evidência coletada indica que nenhum modelo supera consistentemente todos os outros em todos os alvos, justificando a estratégia adotada neste trabalho de comparar um conjunto amplo e fixo de modelos (Tabela~\ref{tbl:modelos-ml}, Capítulo~\ref{chp:capitulo4}) com otimização de hiperparâmetros individual via Optuna\,\cite{akiba2019optuna}, em vez de restringir a comparação a uma única arquitetura \textit{a priori}. O modelo híbrido de \citeauthor{shao2024hybrid}, que combina um modelo de balanço de oxigênio (OBM) com uma rede neural profunda para predição do tempo de sopro, complementa essa perspectiva ao demonstrar que variáveis físico-derivadas específicas do processo elevam o desempenho preditivo\,\cite{shao2024hybrid}.

%-------------------------------------------------------------------%
\section{Modelagem de Fornos EAF e Interpretabilidade}
\label{sec:eaf-interpretabilidade}

Embora o foco desta dissertação seja o BOF, estudos em fornos a arco elétrico (EAF) oferecem contribuições metodológicas relevantes, especialmente no que tange à interpretabilidade de modelos de ML aplicados a processos siderúrgicos. \citeauthor{skrjanc2022eaf} avaliaram a importância de variáveis em modelos preditivos de EAF usando análise de importância por permutação, identificando que peso de sucata e temperatura de corrida dominam as predições de processo\,\cite{skrjanc2022eaf}. O estudo demonstra que análise de importância de atributos permite verificar a coerência física das explicações — por exemplo, confirmando que variáveis de carga metálica dominam os resultados, conforme esperado pela termodinâmica.

\citeauthor{logar2022eaf} modelaram o balanço energético de um EAF industrial para estimar perdas térmicas e identificar oportunidades de melhoria operacional, integrando equações de conservação de energia como estrutura de modelos de ML\,\cite{logar2022eaf}. Complementarmente, \citeauthor{manojlovic2022eaf} aplicaram XGBoost com análise SHAP a dados de um EAF industrial, mostrando que os atributos de maior importância SHAP (uso de oxigênio, peso de corrida) correspondiam às variáveis operacionais que especialistas de processo já identificavam como determinantes\,\cite{manojlovic2022eaf}. Esse último trabalho fundamenta diretamente o uso de SHAP como ferramenta de validação física nesta dissertação: espera-se que atributos derivados de balanços estequiométricos (como demanda teórica de O$_2$ e basicidade da escória) ocupem posições de alta importância nos modelos treinados com a configuração físico-informada.

%-------------------------------------------------------------------%
\section{Geração de Dados Sintéticos para Processos Industriais}
\label{sec:dados-sinteticos}

A escassez de dados rotulados é um obstáculo recorrente em aplicações industriais de ML, onde cada corrida representa um custo elevado e os rótulos (composição química final, temperatura) exigem medições analíticas precisas\,\cite{buggineni2024synthetic}. A geração de dados sintéticos emerge como alternativa para viabilizar o treinamento de modelos quando o volume de dados reais é insuficiente.

\citeauthor{buggineni2024synthetic} propuseram um \textit{framework} sistemático para aumentar dados de manufatura com instâncias sintéticas, validando a abordagem em múltiplos cenários industriais\,\cite{buggineni2024synthetic}. Os autores estabelecem que a validade dos dados sintéticos depende do respeito a restrições físicas do processo durante a geração — princípio central da metodologia desta dissertação, implementado via validação por balanço entálpico com tolerância de 10\%. \citeauthor{melo2024synthetic} exploraram a geração de dados sintéticos especificamente no contexto de gêmeos digitais industriais, discutindo métricas de fidelidade que garantem que os dados gerados preservem as distribuições e correlações do processo real\,\cite{melo2024synthetic}.

A escolha de Monte Carlo com validação por balanço físico como técnica de geração é fundamentada na literatura: métodos baseados em amostragem de distribuições físicas são mais adequados do que abordagens generativas (GANs, VAEs) quando o tamanho amostral alvo é pequeno e as restrições físicas são bem conhecidas\,\cite{buggineni2024synthetic,morales2025synthetic}. \citeauthor{morales2025synthetic} reforçam esse ponto em revisão focada em manutenção preditiva industrial, concluindo que dados sintéticos fisicamente coerentes permitem generalização comparável a dados reais para tarefas de regressão de processo\,\cite{morales2025synthetic}. \citeauthor{giannini2022synthetic} demonstraram a geração de datasets de manufatura via simulação de eventos discretos, mostrando que modelos físicos de baixa fidelidade são suficientes para produzir instâncias de treino úteis\,\cite{giannini2022synthetic}.

%-------------------------------------------------------------------%
\section{Posicionamento frente às Arquiteturas Híbridas}
\label{sec:hibridas}

Conforme as três topologias de integração físico-ML descritas na Seção~\ref{sec:arquiteturas-hibridas} (serial, paralela e informada por atributos), os trabalhos revisados em BOF cobrem as três: \citeauthor{chattopadhyay2022hybrid} adotam a abordagem serial\,\cite{chattopadhyay2022hybrid}, \citeauthor{xia2024dmpinn} a paralela\,\cite{xia2024dmpinn}, \citeauthor{shao2024hybrid} também a paralela, com balanço de O$_2$\,\cite{shao2024hybrid}, e \citeauthor{bae2020bof} a informada por atributos\,\cite{bae2020bof}. Esta dissertação posiciona-se na mesma terceira topologia de \citeauthor{bae2020bof} com a Config~C — a novidade não está, portanto, na topologia em si (Seção~\ref{sec:analise-lacunas} detalha a lacuna real).

%-------------------------------------------------------------------%
\section{Tabela Comparativa e Análise das Lacunas}
\label{sec:analise-lacunas}

Nenhum dos trabalhos revisados realiza simultaneamente as três operações que definem a proposta desta dissertação: (i) isolar o efeito da engenharia de atributos do efeito de capacidade arquitetural mediante um protocolo A×B×C controlado, com o mesmo conjunto de modelos e protocolo de \textit{tuning} em ambos os braços da comparação; (ii) usar SHAP não apenas como ferramenta de ranqueamento de importância, mas como mecanismo de validação de coerência física, testando hipóteses pré-registradas sobre a direção e a posição relativa dos atributos físico-derivados; e (iii) tratar temperatura, carbono e fósforo como três alvos independentes sobre o mesmo protocolo experimental. Como a Tabela~\ref{tbl:trabalhos-relacionados} deixa explícito, a topologia informada por atributos por si só \emph{não} é a lacuna — \citeauthor{bae2020bof} já a aplica a BOF nos três alvos\,\cite{bae2020bof}; a lacuna está no protocolo de avaliação controlado que isola a contribuição específica da engenharia de atributos.

\begin{table}[H]
\centering
\caption{Síntese dos principais trabalhos relacionados a PIML em processos siderúrgicos. N: tamanho amostral; Dados: reais de planta industrial (real) ou sintéticos/simulados (sint.).}
\label{tbl:trabalhos-relacionados}
\scriptsize
\resizebox{\textwidth}{!}{%
\setlength{\tabcolsep}{3pt}
\begin{tabular}{|p{2.6cm}|p{1.1cm}|p{2.3cm}|p{1.5cm}|p{1.3cm}|p{1.7cm}|p{2.3cm}|}
\hline
\textbf{Referência} & \textbf{Processo} & \textbf{Abordagem física} & \textbf{Eng.\ Atrib.} & \textbf{Alvos} & \textbf{N; Dados} & \textbf{Contribuição} \\ \hline
\citeauthor{xia2024dmpinn} (\citeyear{xia2024dmpinn})\,\cite{xia2024dmpinn} & BOF & Paralela: eq.\ cinética + rede & Nenhuma (perda) & C & N/D; real & +5,23\,pp hit rate \\ \hline
Chattopadhyay et al. (\citeyear{chattopadhyay2022hybrid})\,\cite{chattopadhyay2022hybrid} & BOF & Serial: termodin.\ + ANNs & Implícita & T, C, P & N/D; real & $-$9,77\% NRMSE P \\ \hline
\citeauthor{shao2024hybrid} (\citeyear{shao2024hybrid})\,\cite{shao2024hybrid} & BOF & Paralela: balanço O$_2$ + DNN & Balanço O$_2$ & Tempo sopro & N/D; real & Features balanço O$_2$ \\ \hline
\citeauthor{bae2020bof} (\citeyear{bae2020bof})\,\cite{bae2020bof} & BOF & 17 features termod. & Físico-deriv. & T, C, P & N/D; real & $+$2--4\,pp hit rate \\ \hline
\citeauthor{zhou2024bof} (\citeyear{zhou2024bof})\,\cite{zhou2024bof} & BOF & Off-gas temporal & Domínio & T, C & 13.528; real & Benchmark SOTA \\ \hline
\citeauthor{wang2020endpoint} (\citeyear{wang2020endpoint})\,\cite{wang2020endpoint} & BOF & \textit{Review} & --- & C & --- & Síntese de métodos de predição de C \\ \hline
\citeauthor{ghalati2023review} (\citeyear{ghalati2023review})\,\cite{ghalati2023review} & BOF & \textit{Review} & --- & T, C, P & --- & Síntese de ML aplicado a BOF \\ \hline
\citeauthor{skrjanc2022eaf} (\citeyear{skrjanc2022eaf})\,\cite{skrjanc2022eaf} & EAF & Importância variáveis & Nenhuma & Energia & N/D; real & Domin.\ peso sucata \\ \hline
\citeauthor{logar2022eaf} (\citeyear{logar2022eaf})\,\cite{logar2022eaf} & EAF & Balanço energético & Implícita & Energia & N/D; real & Híbrido físico-ML \\ \hline
\citeauthor{manojlovic2022eaf} (\citeyear{manojlovic2022eaf})\,\cite{manojlovic2022eaf} & EAF & XGBoost + SHAP & Nenhuma & Energia & N/D; real & SHAP valida coerência física (O$_2$, peso corrida) \\ \hline
\citeauthor{wang2025phoenix} (\citeyear{wang2025phoenix})\,\cite{wang2025phoenix} & Manuf. (solda) & Físico-informado & Físico-deriv. & Múltiplos & N/D; real & Escalabilidade em manufatura complexa \\ \hline
\textbf{Esta dissertação} & \textbf{BOF} & \textbf{Informada por atributos} & \textbf{Físico-deriv.\ (8)} & \textbf{T, C, P} & \textbf{10.000; sint.} & \textbf{Protocolo A$\times$B$\times$C + SHAP como validação} \\ \hline
\end{tabular}}
\end{table}

A análise da Tabela~\ref{tbl:trabalhos-relacionados} revela dois aspectos relevantes, nenhum deles sendo a topologia em si. Primeiro, a coluna N;\,Dados deixa explícito o contraste central desta dissertação frente à literatura: os trabalhos em BOF revisados usam dados reais de planta (a maioria sem reportar N publicamente), enquanto esta dissertação usa 10.000 corridas sintéticas — uma limitação de validade externa discutida no Capítulo~\ref{chp:con}, não uma vantagem metodológica. Segundo, mesmo entre os trabalhos que usam atributos físico-derivados (\citeauthor{bae2020bof}\,\cite{bae2020bof}, \citeauthor{wang2025phoenix}\,\cite{wang2025phoenix}), nenhum isola o efeito da engenharia de atributos de um efeito de capacidade arquitetural via comparação controlada com o mesmo conjunto de modelos e protocolo de \textit{tuning} em ambos os lados.

A literatura existente raramente avalia o impacto das features físico-derivadas de forma isolada, controlando outros fatores. Comparações entre conjuntos de features (brutas vs. brutas+físicas) com a mesma família de modelos e mesmo protocolo de tuning são escassas, tornando difícil atribuir ganhos de desempenho especificamente à engenharia de atributos\,\cite{zhou2024bof,chattopadhyay2022hybrid}. Esta dissertação busca preencher essa lacuna ao adotar três configurações experimentais comparáveis — Config~A (brutas, conjunto completo de modelos), Config~B (brutas, MLP profunda) e Config~C (brutas + físico-derivadas, conjunto completo) — que permitem isolar o efeito da engenharia de atributos da escolha de arquitetura. Os limiares de qualidade industrial adotados como critério de avaliação são MAE $\leq 10$\,°C para temperatura, $\leq 0{,}020$\,\%~em~massa para carbono e $\leq 0{,}003$\,\%~em~massa para fósforo\,\cite{zhou2024bof,bae2020bof} — valores inspirados nos limiares de \textit{hit rate} reportados nesses trabalhos, mas não diretamente equivalentes a eles: um \textit{hit rate} de $\pm10$\,°C descreve a fração de corridas dentro de uma janela, não um limite superior de MAE agregado, e o MAE efetivamente associado a esse \textit{hit rate} em \citeauthor{zhou2024bof} é de aproximadamente 5,78\,°C (Seção~\ref{sec:bof-endpoint}) — um critério mais rigoroso do que os 10\,°C adotados aqui.

Finalmente, embora \citeauthor{bae2020bof} e \citeauthor{zhou2024bof} já tratem o fósforo como alvo, \citeauthor{xia2024dmpinn} foca exclusivamente em carbono e \citeauthor{chattopadhyay2022hybrid} modela fósforo apenas como parte secundária do pipeline serial. Esta dissertação trata T\textsubscript{final}, C\textsubscript{final} e P\textsubscript{final} como três alvos independentes sob o mesmo protocolo controlado, avaliando se a basicidade da escória — atributo físico-derivado diretamente relacionado à eficiência de desfosforação\,\cite{bae2020bof} — aparece como feature relevante nos modelos via análise SHAP.
